\chapter{Version Control with Git}

\section{Why Use Version Control?}\index{version control}\index{Git}
As you write programs, you change them. A change may break something that
worked before, and you may want to return to an earlier version. \newterm{Git}
records the history of a project so you can compare and undo changes. A project
tracked by Git is a \newterm{repository}.

\section{Creating a Repository}
Make a folder for your project, then run:

\begin{minted}{text}
git init
\end{minted}

This creates a hidden \texttt{.git} folder that stores the history. Nothing is
recorded yet.

\section{Saving a Snapshot: Commits}\index{commit}
A \newterm{commit} is a saved snapshot of the files you chose to track. Create a file
and save it:

\begin{minted}{text}
echo "hello" > notes.txt
git add notes.txt
git commit -m "Add notes"
\end{minted}

\texttt{git add} stages a file, which means it will be included in the next
commit. \texttt{git commit} saves the staged files with a message that
describes the change. Write messages that explain why the change was made.

To see the history, run \texttt{git log --oneline}. Each commit has a short
identifier and its message.

\textbf{Worked example: a first commit.} Start in an empty folder and run \texttt{git init}. Create the file, stage it, and commit it:
\begin{minted}{text}
echo "hello" > notes.txt
git add notes.txt
git commit -m "Add notes"
git log --oneline
\end{minted}
The last command shows one line for each commit. Your line will show a different identifier, but it will show your message: \texttt{Add notes}. The identifier is a short form of the commit's unique hash.

\section{Branches}\index{branch}
A \newterm{branch} is a separate line of work. Making a branch lets you try an
idea without changing the main line of your project:

\begin{minted}{text}
git branch new-idea
git switch new-idea
\end{minted}

Commits made on \texttt{new-idea} do not appear on the main branch. Run
\texttt{git switch} with the name of the main branch to return to it. If the
idea works, you can merge it back in. If it does not, you can delete the branch.

\section{A Routine for Each Session}
\begin{enumerate}
\item Check what changed with \texttt{git status}.
\item Stage the files you want to save with \texttt{git add}.
\item Commit them with a clear message.
\item Repeat after each small step that works.
\end{enumerate}

\begin{Exercise}[label=vc1, title=Record a change]
You edit \texttt{notes.txt} and want to save the change. Which three commands do
you run, in order?
\end{Exercise}
\begin{Answer}[ref=vc1]
\texttt{git status} (optional, to check), then \texttt{git add notes.txt}, then
\texttt{git commit -m "Describe the change"}.
\end{Answer}

\begin{Exercise}[label=vc2, title=Branch or main?]
You want to try a new version of a function without risking the working code.
What should you do?
\end{Exercise}
\begin{Answer}[ref=vc2]
Create a branch with \texttt{git branch} and switch to it with \texttt{git switch}.
Work and commit there. Switch back to the main branch if the experiment fails.
\end{Answer}
